\documentclass[11pt,a4paper]{article}
\usepackage[utf8]{inputenc}
\usepackage[T1]{fontenc}
\usepackage[norsk]{babel}
\usepackage{tgtermes}
\usepackage{tgheros}
\renewcommand{\ttdefault}{txtt}
\usepackage{amsmath,amssymb}
\usepackage[a4paper,left=2.45cm,right=2.6cm,top=2.3cm,bottom=2.2cm,headsep=0.9cm]{geometry}
\usepackage{fancyhdr}
\usepackage{lastpage}
\usepackage{xcolor}
\usepackage{tikz}
\usepackage{pgfplots}
\pgfplotsset{compat=1.18}
\usepackage{booktabs}
\usepackage{array}
\usepackage{enumitem}
\usepackage{listings}
\usepackage{needspace}
\usepackage{multirow}
\usetikzlibrary{positioning,arrows.meta,calc}
\definecolor{insgrey}{gray}{0.42}
\definecolor{radiogrey}{gray}{0.70}
\setlength{\parindent}{0pt}
\setlength{\parskip}{0.45em}
\setlist[itemize]{itemsep=0pt,topsep=2pt,label=\textbullet,font=\sffamily}
\lstset{basicstyle=\ttfamily\small,columns=fullflexible,keepspaces=true,
  showstringspaces=false,upquote=true,aboveskip=0.6em,belowskip=0.6em,
  literate={æ}{{\ae}}1 {ø}{{\o}}1 {å}{{\aa}}1 {Æ}{{\AE}}1 {Ø}{{\O}}1 {Å}{{\AA}}1}
\pagestyle{fancy}
\fancyhf{}
\renewcommand{\headrulewidth}{0pt}
\fancyhead[R]{\fontsize{7.5}{9}\selectfont\sffamily\color{insgrey}TDT4172 - Introduksjon til maskinlæring - Øvingssett 9}
\fancyfoot[R]{\fontsize{7}{9}\selectfont\sffamily\color{insgrey}\thepage/\pageref{LastPage}}
\newcommand{\delstart}[1]{\clearpage
  \noindent\hspace*{-1.55em}\makebox[1.55em][l]{\sffamily\bfseries\fontsize{12}{14}\selectfont\raisebox{0.35em}{#1}}%
  {\sffamily\bfseries\fontsize{18}{22}\selectfont Part #1}\par\vspace{1.0em}}
\newcommand{\sporsmal}[1]{\par\Needspace{9\baselineskip}\vspace{1.9em}%
  \noindent{\bfseries Spørsmål #1}\par\vspace{-0.45em}%
  \noindent{\bfseries -{}-{}-{}-{}-{}-{}-{}-{}-{}-{}-{}-{}-{}-{}-{}-{}-}\par\vspace{-0.15em}}
\newcommand{\velg}[1]{\par\Needspace{6\baselineskip}\noindent{\bfseries #1}\par\vspace{0.3em}}
\newcommand{\radio}[1]{\tikz[baseline=-0.62ex]{\draw[radiogrey,line width=0.45pt,fill=gray!5] (0,0) circle (0.56em);
  \node[inner sep=0pt,font=\sffamily\fontsize{5.5}{6}\selectfont,text=insgrey] at (0,0) {#1};}}
\newenvironment{alternativer}{\begin{list}{}{\setlength{\leftmargin}{2.35em}\setlength{\labelwidth}{1.4em}%
  \setlength{\labelsep}{0.6em}\setlength{\itemsep}{0.85em}\setlength{\topsep}{0.45em}\setlength{\parsep}{0pt}}%
  \sffamily\fontsize{10}{12.6}\selectfont}{\end{list}}
\newcommand{\alt}[2]{\item[\radio{#1}]#2}
\newcommand{\makspoeng}{\par\vspace{2.2em}\hfill{\sffamily\fontsize{8.5}{10}\selectfont Maks poeng: 10}\hspace*{0.5em}\par}
\begin{document}

\noindent\hspace*{-1.55em}\makebox[1.55em][l]{\raisebox{0.3em}{\sffamily\bfseries\fontsize{11}{12}\selectfont i}}%
{\sffamily\bfseries\fontsize{18}{22}\selectfont Frontpage}\par\vspace{1.1em}

{\bfseries Eksamensoppgave i TDT4172 - Introduksjon til maskinlæring}\\
{\bfseries Øvingssett:} 9 av 10 (uoffisielt øvingssett)\\
{\bfseries Tid:} 4 timer

\vspace{0.6em}
{\bfseries Faglig kontakt:} --- (øvingssett, ikke en offisiell eksamen)\\
{\bfseries Møter i eksamenslokalet:} Nei

\vspace{0.6em}
{\bfseries Hjelpemiddelkode/Tillatte hjelpemidler:} D / Ingen trykte eller håndskrevne hjelpemidler tillatt.
Bestemt, enkel kalkulator tillatt.

\vspace{0.6em}
{\bfseries\underline{ANNEN INFORMASJON:}}

\vspace{0.4em}
{\bfseries Les oppgavene nøye og gjør dine egne antagelser.} Presiser i besvarelsen hvilke forutsetninger du
har lagt til grunn i tolkning/avgrensing av oppgaven.

Dette øvingssettet er laget for egentrening, med samme format og vanskelighetsgrad som de
siste ordinære eksamenene i emnet. Løsningsforslaget står på siste side.

\vspace{1.6em}
{\bfseries\underline{FAGSPESIFIKK INFORMASJON}}

\vspace{0.4em}
{\bfseries Ingen papirhåndtegninger:}\\
Denne eksamenen tillater ikke bruk av håndtegninger.

\vspace{1.4em}
{\bfseries Vekting av oppgavene}:\\
Denne eksamenen inneholder flervalgsoppgaver organisert i ti grupper på fire spørsmål hver. Hver
gruppe vurderes separat etter følgende regler:

\vspace{0.6em}
For hver av de fire spørsmålene i gruppen får du følgende poeng
\begin{itemize}[leftmargin=2.2em,label=\textperiodcentered,itemsep=0pt]
\item {\bfseries Riktig svar:} +2,5 poeng
\item {\bfseries Feil svar (et alternativ er valgt, men ikke korrekt):} $-$0,83 poeng
\item {\bfseries Ikke besvart:} 0 poeng
\end{itemize}
{\bfseries Merk} at dersom du først har valgt et svaralternativ, kan du {\bfseries ikke velge det bort} igjen. Da må du
svare på det spørsmålet. Med andre ord mister du muligheten til å la et spørsmål stå blankt så fort
du har trykket på et alternativ.

\vspace{0.6em}
Poengsummen for hver gruppe er {\bfseries summen} av poengene for de fire enkeltspørsmålene.\\
For {\bfseries alle spørsmålene} gjelder en {\bfseries null-grense}: dersom summen for denne gruppen blir negativ,
settes verdien til {\bfseries 0}. Den totale poengsummen på eksamen er summen av poengene fra alle
gruppene.

\delstart{1}
\sporsmal{1}
Et e-postprogram skal lære å sortere innkommende e-post som \emph{spam} eller \emph{ikke spam}.
Programmet trenes på 50\,000 e-poster som brukere har merket, og utviklerne rapporterer
\textbf{accuracy} på et usett testsett.

I Mitchells definisjon lærer et program fra \textbf{erfaring E} med hensyn til en \textbf{oppgave T}
og et \textbf{ytelsesmål P}. Utviklerne vurderer fire endringer.

Hvilken endring er en endring av \textbf{P}, mens E og T er uendret?
\velg{Velg ett alternativ:}
\begin{alternativer}
\alt{A}{Å bytte fra logistisk regresjon til et beslutningstre trent på de samme dataene}
\alt{B}{Å la programmet også fordele e-post i mappene «jobb» og «privat»}
\alt{C}{Å rapportere recall for spam-klassen i stedet for accuracy på testsettet}
\alt{D}{Å trene på 80\,000 merkede e-poster i stedet for på de 50\,000 opprinnelige}
\end{alternativer}
\sporsmal{2}
Du utforsker Titanic-data med 400 passasjerer og plotter et \textbf{histogram} over featuren
\emph{Age}, ett for \emph{Survived}=0 og ett for \emph{Survived}=1. Begge fordelingene ser svært
hakkete ut, med mange tomme stolper innimellom høye enkeltstolper.

Hva er den beste tolkningen av plottet?
\velg{Velg ett alternativ:}
\begin{alternativer}
\alt{A}{Age burde vært vist i et bar-plot, som viser fordelingen av en kontinuerlig variabel per klasse}
\alt{B}{Age burde min-max-skaleres før plotting, ellers blir stolpehøydene i histogrammet ujevne}
\alt{C}{Det er valgt for få bins, slik at hver bin inneholder for mange datapunkter til å vise formen}
\alt{D}{Det er valgt for mange bins, slik at hver bin har for få datapunkter til å vise formen}
\end{alternativer}
\sporsmal{3}
Du bruker gradient descent til å minimere tapsfunksjonen $L(w) = w^2$, som har minimum i $w=0$.
Oppdateringsregelen er $w_{t+1} = w_t - \eta \, \frac{dL}{dw}(w_t)$, startverdien er $w_0 = 2$ og
læringsraten er $\eta = 1.5$.

Hvilken beskrivelse passer med de tre første stegene?
\Needspace{11\baselineskip}
\velg{Velg ett alternativ:}
\begin{alternativer}
\alt{A}{$w$ blir $1,\ 0.5,\ 0.25$: den avtar jevnt mot null, men konvergerer ganske tregt}
\alt{B}{$w$ blir $-1,\ 0.5,\ -0.25$: den skifter fortegn og krymper, så den konvergerer}
\alt{C}{$w$ blir $0,\ 0,\ 0$: den treffer minimum i første steg og blir der etterpå}
\alt{D}{$w$ blir $-4,\ 8,\ -16$: den skifter fortegn og vokser, så treningen divergerer}
\end{alternativer}
\sporsmal{4}
Du trener logistisk regresjon og plotter tap og accuracy på \textbf{treningsdataene} etter hver epoke.
Fra epoke 40 til epoke 200 synker tapet jevnt fra 0.35 til 0.28, mens accuracy ligger helt flatt på 0.86.

Hva er den mest sannsynlige forklaringen?
\velg{Velg ett alternativ:}
\begin{alternativer}
\alt{A}{Modellen har begynt å overtilpasse, siden tapet fortsetter å synke uten at modellen blir bedre på dataene}
\alt{B}{Prediksjonene $\hat y$ flyttes nærmere 0 eller 1 uten at noen krysser terskelen, så klassene blir de samme}
\alt{C}{Læringsraten er så lav at vektene endres for lite til at prediksjonene $\hat y$ påvirkes i det hele tatt}
\alt{D}{Tapet synker bare fordi bias $b$ endres, og en endring i $b$ alene kan ikke endre noen klassifisering}
\end{alternativer}
\makspoeng

\delstart{2}
\sporsmal{1}
En modell trenes på et datasett der 99\,\% av pasientene er friske og 1\,\% er syke. Den får 99\,\%
accuracy, men predikerer nesten aldri at en pasient er syk, selv for pasienter med tydelige symptomer.

Hvilken forklaring er riktig?
\velg{Velg ett alternativ:}
\begin{alternativer}
\alt{A}{Prediksjonen domineres av den skjeve klassefordelingen, så informasjonen i featurene veier for lite}
\alt{B}{De skjeve klassene gjør tapsflaten ikke-konveks, så gradient descent stopper i et lokalt minimum}
\alt{C}{Featurene blir avhengige av hverandre når klassene er skjeve, så modellens antakelser brytes}
\alt{D}{Modellen overtilpasser til de syke, siden hvert sykt punkt får uforholdsmessig stor vekt i tapet}
\end{alternativer}
\sporsmal{2}
Et testsett har 1\,000 punkter, hvorav 50 er positive. En defekt modell gir \emph{samme} skår $0.3$ til
alle punktene, og prediksjonen er positiv hvis skåren er over $0.5$.

Hva blir accuracy og ROC AUC for modellen?
\velg{Velg ett alternativ:}
\begin{alternativer}
\alt{A}{Accuracy $0.95$ og ROC AUC $0.50$}
\alt{B}{Accuracy $0.05$ og ROC AUC $0.00$}
\alt{C}{Accuracy $0.05$ og ROC AUC $0.50$}
\alt{D}{Accuracy $0.95$ og ROC AUC $0.00$}
\end{alternativer}
\sporsmal{3}
Et datasett for en sjelden sykdom har 4\,\% positive. En «modell» gir hvert punkt en tilfeldig skår
mellom 0 og 1, uavhengig av featurene. Du tegner \textbf{precision-recall-kurven} ved å variere
terskelen.

Hvordan ser kurven omtrent ut?
\velg{Velg ett alternativ:}
\begin{alternativer}
\alt{A}{Precision faller lineært fra 1.00 til 0.00 når recall øker fra 0 til 1, som diagonalen i ROC-rommet}
\alt{B}{Precision øker fra 0.04 mot 1.00 når recall øker, siden en lavere terskel finner flere av de positive}
\alt{C}{Precision ligger rundt 0.50 for alle recall-verdier, siden en tilfeldig modell treffer riktig i omtrent halvparten av tilfellene}
\alt{D}{Precision ligger rundt 0.04 for alle recall-verdier, fordi andelen positive blant de flaggede er som i dataene}
\end{alternativer}
\sporsmal{4}
En klassifikator gir følgende resultat på et testsett: TP $= 45$, FP $= 5$, FN $= 45$ og TN $= 105$.

Hva er \textbf{F1-scoren}?
\velg{Velg ett alternativ:}
\begin{alternativer}
\alt{A}{$0.75$}
\alt{B}{$0.70$}
\alt{C}{$0.90$}
\alt{D}{$0.64$}
\end{alternativer}
\makspoeng

\delstart{3}
\sporsmal{1}
En node i et klassifikasjonstre har 8 positive og 2 negative datapunkter og predikerer positiv klasse.
En kandidatsplitt gir venstre barn med 4 positive og 0 negative, og høyre barn med 4 positive og
2 negative. Begge barna predikerer dermed fortsatt positiv klasse, og feilraten er 0.20 både før og etter.
Treet bruker Gini $1-\sum_i p_i^2$, vektet med andelen punkter i hvert barn.

Hvilket utsagn er riktig?
\velg{Velg ett alternativ:}
\begin{alternativer}
\alt{A}{Gini synker fra 0.32 til 0.22, som er snittet av Gini-urenheten i de to barnenodene}
\alt{B}{Vektet Gini synker fra 0.32 til 0.27, så splitten kan velges selv om ingen prediksjon endres}
\alt{C}{Splitten gir ingen reduksjon i urenhet, siden andelen feilklassifiserte punkter er 0.20 både før og etter}
\alt{D}{Vektet Gini er uendret, siden begge barna predikerer samme klasse som foreldrenoden}
\end{alternativer}
\sporsmal{2}
Featuren \emph{Region} har verdiene Midt, Nord, Sør og Vest, uten noen naturlig rekkefølge. Den kodes med
\texttt{LabelEncoder}, som gir Midt $=0$, Nord $=1$, Sør $=2$ og Vest $=3$. Et beslutningstre splitter på den
kodede featuren med regler på formen $\text{Region} \le t$.

Hvilken todeling av regionene kan treet \textbf{ikke} oppnå med én enkelt splitt?
\velg{Velg ett alternativ:}
\begin{alternativer}
\alt{A}{$\{\text{Midt}, \text{Sør}\}$ mot $\{\text{Nord}, \text{Vest}\}$}
\alt{B}{$\{\text{Midt}, \text{Nord}, \text{Sør}\}$ mot $\{\text{Vest}\}$}
\alt{C}{$\{\text{Midt}\}$ mot $\{\text{Nord}, \text{Sør}, \text{Vest}\}$}
\alt{D}{$\{\text{Midt}, \text{Nord}\}$ mot $\{\text{Sør}, \text{Vest}\}$}
\end{alternativer}
\sporsmal{3}
En node har 8 positive og 8 negative datapunkter. En splitt gir venstre barn med 6 positive og
0 negative, og høyre barn med 2 positive og 8 negative.

Hva er \textbf{informasjonsgevinsten} $\Delta I = I_\text{før} - E[I_\text{etter}]$ for denne splitten,
med entropi beregnet med $\log_2$?
\velg{Velg ett alternativ:}
\begin{alternativer}
\alt{A}{$0.45$}
\alt{B}{$0.55$}
\alt{C}{$0.64$}
\alt{D}{$0.38$}
\end{alternativer}
\sporsmal{4}
Titanic-dataene har featuren \emph{PassengerId}, et unikt nummer for hver passasjer. Treet lager én gren
per verdi av en kategorisk feature, som i pseudokoden fra forelesningen, og velger i hver node featuren
med størst informasjonsgevinst. \emph{PassengerId} er tatt med blant featurene.

Hva skjer i rotnoden?
\velg{Velg ett alternativ:}
\begin{alternativer}
\alt{A}{\emph{PassengerId} gir høy gevinst og generaliserer godt, siden nye passasjerer følger grenen med nærmeste nummer}
\alt{B}{\emph{PassengerId} gir gevinst 0, siden nummeret ikke henger sammen med overlevelse, og blir derfor aldri valgt}
\alt{C}{\emph{PassengerId} gir størst mulig gevinst, siden hver gren får ett punkt, men splitten er ubrukelig på nye passasjerer}
\alt{D}{\emph{PassengerId} gir negativ gevinst, siden $E[I_\text{etter}]$ øker med antall grener, og blir derfor ikke valgt}
\end{alternativer}
\makspoeng

\delstart{4}
\sporsmal{1}
Du tilpasser en \textbf{lineær} regresjonsmodell $f(x)=\beta_0+\beta_1 x$ til data som tydelig følger
en parabel. Både trenings- og valideringstapet flater ut på et høyt nivå, og de to kurvene ligger
nesten oppå hverandre. En kollega foreslår å samle inn dobbelt så mye data av samme type.

Hva er det mest sannsynlige utfallet av å doble datamengden?
\velg{Velg ett alternativ:}
\begin{alternativer}
\alt{A}{Lite endring, fordi feilen domineres av bias som mer data av samme type ikke fjerner}
\alt{B}{Treningstapet går klart ned, fordi modellen får flere punkter å tilpasse linjen til}
\alt{C}{Valideringstapet går klart ned, fordi mer data reduserer både bias og varians}
\alt{D}{Valideringstapet øker, fordi flere punkter gir mer støy som modellen overtilpasser til}
\end{alternativer}
\sporsmal{2}
Du har tre features $x_1, x_2, x_3$ og lager polynomfeatures av grad 2, med klassen
\texttt{PolynomialFeatures} og \texttt{degree=2}. Utdataene inneholder konstantkolonnen, de opprinnelige
featurene, alle kvadrater og alle produkter av to ulike features.

Hvor mange kolonner får den transformerte datamatrisen?
\velg{Velg ett alternativ:}
\begin{alternativer}
\alt{A}{$10$}
\alt{B}{$9$}
\alt{C}{$7$}
\alt{D}{$6$}
\end{alternativer}
\sporsmal{3}
Dataene er generert som $y = \lfloor x \rfloor$ pluss litt støy, for $x$ jevnt fordelt i $[0, 4)$, og
delt tilfeldig i trening og test. Et polynom av grad 5 tilpasset med lineær regresjon bølger seg rundt
trappetrinnene, mens et \textbf{regresjonstre} med dybde 2 får klart lavere test-MSE.

Hva er den beste forklaringen?
\velg{Velg ett alternativ:}
\begin{alternativer}
\alt{A}{Treet kan følge trinnene videre utenfor treningsområdet, mens polynomet bare er pålitelig innenfor dette området}
\alt{B}{Treet gir én konstant verdi per intervall av $x$, mens polynomet er glatt og ikke kan gjøre brå sprang}
\alt{C}{Polynomet ville passet like godt om $x$ var standardisert først, siden høye potenser av $x$ gir dårlig tilpasning}
\alt{D}{Polynomet har for mange parametre for dataene, mens treet med dybde 2 bare har to parametre som tilpasses}
\end{alternativer}
\sporsmal{4}
Du trener en lineær regresjonsmodell med $\text{MSE} = \frac{1}{2N}\sum_i \bigl(y^{(i)} - f(x^{(i)})\bigr)^2$.
En kollega bruker samme data, startverdier og læringsrate, men definerer MSE uten faktoren $\tfrac12$.

Hvilket utsagn er riktig?
\velg{Velg ett alternativ:}
\begin{alternativer}
\alt{A}{Begge har minimum i de samme parameterverdiene, men kollegaens gradient er dobbelt så stor, som en dobbel læringsrate}
\alt{B}{Begge gir nøyaktig de samme oppdateringene, siden faktoren $\tfrac12$ forsvinner når tapet deriveres}
\alt{C}{Kollegaens tap er ikke lenger konvekst, siden faktoren $\tfrac12$ trengs for at MSE skal ha ett minimum}
\alt{D}{Kollegaens minimum ligger ved parameterverdier som er dobbelt så store, siden tapet er dobbelt så stort overalt}
\end{alternativer}
\makspoeng

\delstart{5}
\sporsmal{1}
Du har målt responstiden til en server 200 ganger og beregnet \textbf{medianen}. Du har ikke
mulighet til å samle inn flere målinger, men vil anslå hvor usikker medianen er.

Hvilken fremgangsmåte er bootstrapping?
\velg{Velg ett alternativ:}
\begin{alternativer}
\alt{A}{Trekk mange utvalg på 100 uten tilbakelegging, beregn medianen i hvert, og se på spredningen}
\alt{B}{Del målingene i 10 disjunkte deler, beregn medianen i hver del, og bruk gjennomsnittet av dem}
\alt{C}{Trekk mange utvalg på 200 med tilbakelegging, beregn medianen i hvert, og se på spredningen}
\alt{D}{Trekk mange utvalg på 200 med tilbakelegging, og bruk medianen av alle trekningene samlet}
\end{alternativer}
\sporsmal{2}
Datapunktene er jevnt fordelt i de fire kvadrantene av $(x_1, x_2)$-planet, like mange i hver.
Punktene i kvadranten der både $x_1 > 0$ og $x_2 > 0$ har klasse 1, alle andre har klasse 0.
Du trener en \textbf{trestump} (én rotnode med to løvnoder) på dataene.

Hva er den høyeste accuracy trestumpen kan oppnå på treningsdataene?
\velg{Velg ett alternativ:}
\begin{alternativer}
\alt{A}{100\,\%, siden rotnoden kan teste $x_1>0$ og $x_2>0$ i samme splitt}
\alt{B}{50\,\%, siden en svak lærer per definisjon ikke er bedre enn å gjette}
\alt{C}{87.5\,\%, siden én løvnode blir ren og den andre får 75\,\% riktig}
\alt{D}{75\,\%, siden én splitt på bare én feature ikke skiller ut kvadranten}
\end{alternativer}
\sporsmal{3}
Du trener gradient boosting for regresjon med kvadratisk tap $\tfrac12(y-F)^2$. Startmodellen
$F_0$ predikerer gjennomsnittet av targets for alle punkter. Treningsdataene har targets
$y = (4,\ 9,\ 14,\ 21)$.

Første tre trenes på \textbf{pseudo-residualene}, altså den negative gradienten av tapet med hensyn
på $F$. Hvilke verdier trener det første treet på?
\velg{Velg ett alternativ:}
\begin{alternativer}
\alt{A}{$(-7.5,\ -2.5,\ 2.5,\ 9.5)$}
\alt{B}{$(64,\ 9,\ 4,\ 81)$}
\alt{C}{$(8,\ 3,\ -2,\ -9)$}
\alt{D}{$(-8,\ -3,\ 2,\ 9)$}
\end{alternativer}
\sporsmal{4}
Tre regresjonsmodeller i et parallelt ensemble predikerer prisen på en bolig med sann pris 5.0 (millioner kr).
Prediksjonene er 5.6, 4.4 og 5.3, og ensemblet predikerer gjennomsnittet av de tre.

Hva er kvadratfeilen til ensemblets prediksjon, og hva er gjennomsnittet av de tre modellenes kvadratfeil?
\velg{Velg ett alternativ:}
\begin{alternativer}
\alt{A}{Ensemblet: $0.01$, snittet av modellene: $0.50$}
\alt{B}{Ensemblet: $0.27$, snittet av modellene: $0.27$}
\alt{C}{Ensemblet: $0.01$, snittet av modellene: $0.27$}
\alt{D}{Ensemblet: $0.10$, snittet av modellene: $0.50$}
\end{alternativer}
\makspoeng

\delstart{6}
\sporsmal{1}
Et fullt koblet lag tar inn aktiveringene $\mathbf{a}^{(l-1)}$ fra et lag med 4 noder og har selv 6 noder.
Aktiveringene skrives på matriseform som $\mathbf{a}^{(l)} = g\bigl(W\mathbf{a}^{(l-1)} + \mathbf{b}\bigr)$,
der $\mathbf{a}^{(l-1)}$ er en kolonnevektor.

Hvilke dimensjoner har $W$ og $\mathbf{b}$?
\velg{Velg ett alternativ:}
\begin{alternativer}
\alt{A}{$W$ er $6 \times 4$ og $\mathbf{b}$ har 4 elementer}
\alt{B}{$W$ er $4 \times 6$ og $\mathbf{b}$ har 4 elementer}
\alt{C}{$W$ er $4 \times 6$ og $\mathbf{b}$ har 6 elementer}
\alt{D}{$W$ er $6 \times 4$ og $\mathbf{b}$ har 6 elementer}
\end{alternativer}
\sporsmal{2}
Et fully connected feed-forward-nettverk (MLP) har et input-lag, to skjulte lag og et output-lag.

Hvilket utsagn om arkitekturen er riktig?
\velg{Velg ett alternativ:}
\begin{alternativer}
\alt{A}{Hver node i et skjult lag får input fra alle nodene i alle tidligere lag, også direkte fra input-laget}
\alt{B}{Nodene i input-laget har egne vekter, bias og aktiveringsfunksjon, akkurat som nodene i de skjulte lagene}
\alt{C}{Hvert skjult lag må ha like mange noder som det er input-features, siden hver node behandler én feature}
\alt{D}{Hver node i et skjult lag får input fra alle nodene i laget rett før, og fra ingen andre noder}
\end{alternativer}
\sporsmal{3}
Et nettverk med én node har $z = wx + b$, $\hat y = \sigma(z)$ og tap $L = \tfrac12(\hat y - y)^2$.
For datapunktet $x = 2$, $y = 1$ er $w = 0.5$ og $b = -1$. Husk at $\sigma'(z) = \sigma(z)(1-\sigma(z))$.

Hva er $\dfrac{\partial L}{\partial w}$ beregnet med kjerneregelen?
\velg{Velg ett alternativ:}
\begin{alternativer}
\alt{A}{$0.250$}
\alt{B}{$-1.000$}
\alt{C}{$-0.250$}
\alt{D}{$-0.125$}
\end{alternativer}
\sporsmal{4}
Et binært klassifiseringsnettverk har to output-noder med softmax. For et datapunkt er logitene
$(z_1, z_2) = (1.5,\ 0.5)$. En kollega vil i stedet bruke én output-node med sigmoid $\sigma(z)$.

Hvilken verdi av $z$ gir samme predikerte sannsynlighet for klasse 1 som softmax-nettverket?
\velg{Velg ett alternativ:}
\begin{alternativer}
\alt{A}{$z = 1.0$}
\alt{B}{$z = 0.5$}
\alt{C}{$z = 1.5$}
\alt{D}{$z = 2.0$}
\end{alternativer}
\makspoeng

\delstart{7}
\sporsmal{1}
Du trekker tilfeldige punkter uniformt fra enhetskuben $[0,1]^d$ og måler euklidsk avstand mellom
par av punkter. Tabellen viser gjennomsnitt og standardavvik for avstandene.
\begin{center}
\begin{tabular}{lccc}
\toprule
$d$ & 2 & 100 & 10\,000 \\
\midrule
Gjennomsnittlig avstand & 0.52 & 4.08 & 40.8 \\
Standardavvik & 0.25 & 0.24 & 0.24 \\
\bottomrule
\end{tabular}
\end{center}
Hvilken konklusjon om clustering i høye dimensjoner støttes best av tabellen?
\velg{Velg ett alternativ:}
\begin{alternativer}
\alt{A}{Avstanden mellom punkter blir stor, så ulike klynger blir lettere å skille enn i to dimensjoner}
\alt{B}{Relativ forskjell mellom nære og fjerne punkter krymper, så avstand skiller dårlig mellom klynger}
\alt{C}{Problemet skyldes for få punkter, så med mange flere datapunkter blir avstandene igjen informative}
\alt{D}{Spredningen i avstand er nesten konstant, så k-means fungerer omtrent like godt i 10\,000 som i to dimensjoner}
\end{alternativer}
\sporsmal{2}
I et k-means-steg er sentroidene $\mathbf{c}_1 = (4,\,4)$, $\mathbf{c}_2 = (3,\,2)$ og
$\mathbf{c}_3 = (0,\,1)$. Punktet $\mathbf{p} = (1,\,4)$ skal tilordnes en klynge med
\textbf{euklidsk avstand}.

Hvilken klynge tilordnes punktet, og hva er avstanden?
\velg{Velg ett alternativ:}
\begin{alternativer}
\alt{A}{Klynge 2, med avstand $8.00$}
\alt{B}{Klynge 3, med avstand $3.16$}
\alt{C}{Klynge 2, med avstand $2.83$}
\alt{D}{Klynge 1, med avstand $3.00$}
\end{alternativer}
\sporsmal{3}
Et datasett har tre klynger og én uteligger:
\begin{center}
\begin{tabular}{ll}
\toprule
Klynge A & $(0,0),\ (1,0),\ (1,1)$ \\
Klynge B & $(6,0),\ (7,0),\ (6,1)$ \\
Klynge C & $(3,6),\ (3,7),\ (4,6)$ \\
Uteligger & $(12,12)$ \\
\bottomrule
\end{tabular}
\end{center}
Du initialiserer k-means med $k=3$ ved å plassere hver nye sentroide på datapunktet som har
\textbf{størst avstand til sin nærmeste} allerede valgte sentroide (euklidsk avstand). Første
sentroide er $(0,0)$, og andre sentroide blir derfor uteliggeren $(12,12)$.

Hvilket punkt blir tredje sentroide?
\velg{Velg ett alternativ:}
\begin{alternativer}
\alt{A}{$(4,\,6)$}
\alt{B}{$(1,\,0)$}
\alt{C}{$(3,\,7)$}
\alt{D}{$(7,\,0)$}
\end{alternativer}
\sporsmal{4}
Du kjører koden under to ganger på nøyaktig de samme dataene \texttt{X}, i samme rekkefølge, uten å
sette \texttt{random\_state}:
\begin{lstlisting}[language=Python]
km = KMeans(n_clusters=3, init="random", n_init=1).fit(X)
db = DBSCAN(eps=0.5, min_samples=5).fit(X)
\end{lstlisting}
Hvilket utsagn om de to kjøringene er riktig?
\velg{Velg ett alternativ:}
\begin{alternativer}
\alt{A}{k-means kan gi ulike inndelinger, mens DBSCAN gir det samme resultatet begge gangene}
\alt{B}{Begge gir det samme resultatet, siden dataene og hyperparametrene er de samme}
\alt{C}{DBSCAN kan gi ulike inndelinger, mens k-means gir det samme resultatet når $k$ er fast}
\alt{D}{Begge kan gi ulike inndelinger, siden begge starter fra tilfeldig valgte punkter}
\end{alternativer}
\makspoeng

\delstart{8}
\sporsmal{1}
Du skal finne anomale målinger i sensordata og sammenligner de tre metodene for anomalideteksjon
fra kurset: k-means, DBSCAN og Isolation Forest.

Hvilken beskrivelse av hvordan hver metode avgjør at et punkt er en anomali, er riktig?
\velg{Velg ett alternativ:}
\begin{alternativer}
\alt{A}{k-means: få splitter til isolasjon; DBSCAN: en avstandsterskel du velger selv; Isolation Forest: label $-1$ fra klyngingen}
\alt{B}{k-means: en avstandsterskel du velger selv; DBSCAN: label $-1$ fra klyngingen; Isolation Forest: få splitter til isolasjon}
\alt{C}{k-means: label $-1$ fra klyngingen; DBSCAN: få splitter til isolasjon; Isolation Forest: en avstandsterskel du velger selv}
\alt{D}{k-means: label $-1$ fra klyngingen; DBSCAN: en avstandsterskel du velger selv; Isolation Forest: få splitter til isolasjon}
\end{alternativer}
\sporsmal{2}
Et sentrert datasett med fire punkter projiseres på første hovedkomponent, og koordinatene langs
PC1 blir $-3,\ -1,\ 1$ og $3$. Kovariansmatrisen beregnes med $N-1$ i nevneren.

Hva er egenverdien $\lambda_1$ som hører til PC1?
\velg{Velg ett alternativ:}
\begin{alternativer}
\alt{A}{$5.00$}
\alt{B}{$20.00$}
\alt{C}{$6.67$}
\alt{D}{$2.58$}
\end{alternativer}
\sporsmal{3}
I SNE er vekten for at $\mathbf{x}_j$ er nabo til $\mathbf{x}_i$ proporsjonal med
$\exp\!\left(-\lVert\mathbf{x}_i-\mathbf{x}_j\rVert^2 / (2\sigma_i^2)\right)$, og normaliseringen er lik
for alle naboer til $\mathbf{x}_i$. Rundt $\mathbf{x}_i$ er $\sigma_i = 1$. Punktet $\mathbf{x}_j$ ligger
i avstand 1 og punktet $\mathbf{x}_k$ i avstand 3 fra $\mathbf{x}_i$.

Omtrent hvor mange ganger større er $p_{j|i}$ enn $p_{k|i}$?
\Needspace{9\baselineskip}
\velg{Velg ett alternativ:}
\begin{alternativer}
\alt{A}{$2.7$}
\alt{B}{$54.6$}
\alt{C}{$9.0$}
\alt{D}{$3.0$}
\end{alternativer}
\sporsmal{4}
To ulike PCA-implementasjoner tilpasses de samme sentrerte dataene. Den ene gir første
hovedkomponent $(0.6,\ 0.8)$, den andre gir $(-0.6,\ -0.8)$. Egenverdiene er like i begge.

Hvilket utsagn er riktig?
\velg{Velg ett alternativ:}
\begin{alternativer}
\alt{A}{Begge er riktige, siden $\mathbf{e}$ og $-\mathbf{e}$ er egenvektorer for samme egenverdi og spenner ut samme akse}
\alt{B}{Bare den første er riktig, siden en hovedkomponent alltid skal ha positive koordinater i originalrommet}
\alt{C}{Projeksjonene på PC1 får motsatt fortegn, så rekonstruksjonen av dataene fra PC1 blir speilvendt i den ene}
\alt{D}{Begge er riktige, men bare fordi dataene er sentrert; for usentrerte data ville fortegnet vært entydig}
\end{alternativer}
\makspoeng

\delstart{9}
\sporsmal{1}
En DQN-agent trenes i FrozenLake. Fra episode 800 til 1500 når den målet i over 90\,\% av episodene.
Rundt episode 1500 faller andelen brått til under 10\,\%, stiger igjen etter noen hundre episoder og
faller så på nytt. I hele perioden er $\epsilon = 0.01$ og læringsraten $0.0005$.

Hva er den mest sannsynlige forklaringen?
\velg{Velg ett alternativ:}
\begin{alternativer}
\alt{A}{Catastrophic forgetting: nettverket overtilpasser til nye overganger; replay buffer og target network hjelper}
\alt{B}{Undertilpasning: Q-nettverket har for få parametre og trenger flere skjulte lag for å representere policyen}
\alt{C}{Reward hacking: agenten har funnet en utilsiktet snarvei i belønningen og utnytter den i stedet for å nå målet}
\alt{D}{For mye exploration: $\epsilon$ er fortsatt for høy, så agenten tar for mange tilfeldige handlinger og bør få $\epsilon = 0$}
\end{alternativer}
\sporsmal{2}
Q-nettverket i en DQN gir for tilstanden $s$ Q-verdiene $(1.2,\ 0.5,\ 2.0,\ -0.3)$ for handlingene
$0,1,2,3$. Agenten utførte handling $1$, fikk reward $r = 0$ og havnet i en ikke-terminal tilstand
$s'$, der target network gir $\max_a Q(s', a) = 1.5$. Diskonteringsfaktoren er $\gamma = 0.9$.

Hva blir MSE-lossen for dette ene datapunktet, slik den beregnes i DQN?
\velg{Velg ett alternativ:}
\begin{alternativer}
\alt{A}{$1.00$}
\alt{B}{$0.97$}
\alt{C}{$3.89$}
\alt{D}{$0.72$}
\end{alternativer}
\sporsmal{3}
I CartPole får agenten reward $+1$ for hvert tidssteg stangen holdes oppe. Anta at agenten aldri mister
stangen, slik at episoden i prinsippet aldri slutter, og at $\gamma = 0.9$.

Hva blir den diskonterte returen $G_t = \sum_{k=0}^{\infty} \gamma^k R_{t+k+1}$ fra et vilkårlig tidspunkt $t$?
\velg{Velg ett alternativ:}
\begin{alternativer}
\alt{A}{$9$}
\alt{B}{$10$}
\alt{C}{$\infty$}
\alt{D}{$1.9$}
\end{alternativer}
\sporsmal{4}
En Q-agent er i tilstand $S_t$ med $Q(S_t, A_t) = 0.4$. Handlingen $A_t$ fører til målet, agenten får
$R_{t+1} = 1$, og miljøet returnerer \texttt{done=True}. Q-tabellen har tilfeldigvis
$\max_A Q(S_{t+1}, A) = 0.8$ for måltilstanden. Bruk
$Q_{new}(S_t,A_t) \leftarrow (1-\alpha)Q(S_t,A_t) + \alpha\,(R_{t+1} + \gamma\max_A Q(S_{t+1},A))$
med $\alpha = 0.2$ og $\gamma = 0.9$, og håndter terminaltilstanden som i forelesningskoden.

Hva blir $Q_{new}(S_t, A_t)$?
\velg{Velg ett alternativ:}
\begin{alternativer}
\alt{A}{$0.66$}
\alt{B}{$0.52$}
\alt{C}{$0.88$}
\alt{D}{$0.60$}
\end{alternativer}
\makspoeng

\delstart{10}
\sporsmal{1}
Du skal forklare tre modeller med SHAP:
(i) en XGBoost-modell for kredittvurdering,
(ii) et konvolusjonsnett for bildeklassifisering som du har trent selv i PyTorch, og
(iii) en ekstern modell som du bare kan sende input til og få prediksjoner tilbake fra via et API.

Hvilken tilordning av explainers til (i), (ii) og (iii) er mulig og mest hensiktsmessig?
\velg{Velg ett alternativ:}
\begin{alternativer}
\alt{A}{(i) TreeExplainer, (ii) DeepExplainer, (iii) DeepExplainer}
\alt{B}{(i) KernelExplainer, (ii) DeepExplainer, (iii) DeepExplainer}
\alt{C}{(i) KernelExplainer, (ii) DeepExplainer, (iii) KernelExplainer}
\alt{D}{(i) TreeExplainer, (ii) DeepExplainer, (iii) KernelExplainer}
\end{alternativer}
\sporsmal{2}
Du beregner SHAP-verdier for 1\,000 testpunkter og lager et \textbf{beeswarm}-plott. Featuren
\emph{Alder} har stor spredning i SHAP-verdiene, mens \emph{Region} nesten bare har verdier nær 0.

Hvilket utsagn om plottet er riktig?
\velg{Velg ett alternativ:}
\begin{alternativer}
\alt{A}{Det er en global forklaring som gjelder for alle mulige input, siden SHAP oppfyller effektivitetsaksiomet}
\alt{B}{Det er satt sammen av lokale forklaringer og beskriver modellen i disse 1\,000 punktene}
\alt{C}{Alder står øverst fordi den har størst gjennomsnittlig SHAP-verdi med fortegn over de 1\,000 punktene}
\alt{D}{Region kan fjernes uten at noen prediksjon endres, siden nullspilleraksiomet gir SHAP-verdier nær 0}
\end{alternativer}
\sporsmal{3}
En student implementerer entropien $I = -\sum_i p_i \log_2 p_i$ for en node i et beslutningstre:
\begin{lstlisting}[language=Python]
import numpy as np

def entropy(counts):
    p = np.array(counts) / np.sum(counts)
    return -np.sum(p * np.log2(p))

print(f"{entropy([3, 1]):.2f} {entropy([4, 0]):.2f}")
\end{lstlisting}
Hva skrives ut?
\velg{Velg ett alternativ:}
\begin{alternativer}
\alt{A}{\texttt{0.81 inf}}
\alt{B}{\texttt{0.81 0.00}}
\alt{C}{\texttt{0.81 nan}}
\alt{D}{\texttt{nan nan}}
\end{alternativer}
\sporsmal{4}
En agent velger handlinger med funksjonen under.
\par\noindent\begin{minipage}{\linewidth}
\begin{lstlisting}[language=Python]
import random
import numpy as np

def select_action(q_values, epsilon):
    if random.random() > epsilon:
        return random.randrange(len(q_values))
    return int(np.argmax(q_values))
\end{lstlisting}
\end{minipage}
Med \texttt{epsilon = 0.05} velger agenten nesten bare tilfeldige handlinger. Hvilken endring gir
riktig $\epsilon$-greedy?
\velg{Velg ett alternativ:}
\begin{alternativer}
\alt{A}{Bytte \texttt{np.argmax} med \texttt{np.argmin} i siste linje}
\alt{B}{Bytte \texttt{random.random()} med \texttt{random.randrange(len(q\_values))}}
\alt{C}{Endre if-testen til \texttt{if random.random() < epsilon:}}
\alt{D}{Returnere \texttt{int(np.argmax(q\_values))} i begge grenene}
\end{alternativer}
\makspoeng

\clearpage
{\sffamily\bfseries\fontsize{18}{22}\selectfont Løsningsforslag}\par\vspace{0.3em}
\noindent\rule{\textwidth}{0.4pt}\par\vspace{1.2em}
Riktig alternativ per spørsmål (bokstaven står i rundingen foran hvert alternativ).\par\vspace{1em}
\begin{center}\sffamily\renewcommand{\arraystretch}{1.35}
\begin{tabular}{l*{4}{>{\centering\arraybackslash}p{1.6cm}}}
\toprule
& \textbf{Sp.\ 1} & \textbf{Sp.\ 2} & \textbf{Sp.\ 3} & \textbf{Sp.\ 4}\\
\midrule
Part 1 & C & D & D & B\\
Part 2 & A & A & D & D\\
Part 3 & B & A & B & C\\
Part 4 & A & A & B & A\\
Part 5 & C & D & D & C\\
Part 6 & D & D & C & A\\
Part 7 & B & C & C & A\\
Part 8 & B & C & B & A\\
Part 9 & A & D & B & B\\
Part 10 & D & B & C & C\\
\bottomrule
\end{tabular}
\end{center}
\end{document}
